\documentclass[french,11pt]{book}
\usepackage{teach}
\usepackage[french]{babel}
\usepackage{amsmath,amssymb}
\usepackage{array}
\newcommand{\R}{\mathbb{R}}
\begin{document}
\begin{center}
{\Large \textbf{Exercices -- Suites numériques}}\\[2mm]
Terminale technologique
\end{center}

\section*{Exercices}
\begin{enumerate}
  \item On considère la suite arithmétique $(u_n)$ définie par $u_0=7$ et $u_{n+1}=u_n+3$. Calculer $u_1$, $u_2$, $u_5$, puis donner une expression de $u_n$ en fonction de $n$.
  \item Une population de bactéries augmente de 12\% chaque heure. Elle vaut $2500$ au départ. Modéliser par une suite, puis estimer la population après 6 heures.
  \item La suite $(v_n)$ est géométrique de premier terme $v_0=80$ et de raison $0,85$. Calculer $v_3$ et interpréter le résultat dans le cas d'une baisse de 15\%.
  \item Un abonnement coûte 18 euros le premier mois puis augmente de 1,50 euro chaque mois. Déterminer le coût du 12e mois et le coût total de la première année.
\end{enumerate}

\end{document}
